\section{Introduction}
% 背景介绍，点出动态图处理的重要性
% 已有工作怎么做的，存在什么问题，面临什么挑战
% 为了解决上述挑战，本文怎么做，实验结果展开说
% 对应贡献与论文组织（重点是贡献，篇幅不够可以不写论文组织结构）
% 先整了个初稿，后面再进一步基于里面的核心优化，做进一步的调整
% Background and motivation
Streaming graph processing supports applications such as fraud detection and network monitoring, where evolving relationships demand timely analytics. Unlike static graph processing, it must maintain graph topology and update analytical results as edges are inserted and deleted. Recomputing from scratch after each update batch repeats substantial work on unchanged regions. Incremental processing reduces this redundancy by reusing existing results and propagating the effects of updates. GPUs offer high parallelism and memory bandwidth for such propagation, yet maintaining dynamic topology on GPUs requires balancing efficient concurrent updates with fast graph traversal. Moreover, large graphs often exceed device memory, necessitating efficient coordination between host-resident topology and GPU computation.

% Existing approaches and limitations
Existing systems address different aspects of this problem. CPU-based systems such as GraphBolt~\cite{2019-GraphBolt} track dependencies across iterations to reuse unaffected computation. GPU-resident systems such as SHGraph~\cite{2020-SHGraph} support parallel topology updates through dynamic graph representations, but require the graph to fit in device memory. CPU--GPU cooperative systems such as Grapin~\cite{2025-Grapin} extend incremental processing to out-of-memory graphs by maintaining topology in host memory and combining GPU hot-subgraph caching with optimized zero-copy access. These approaches provide important building blocks, yet efficient out-of-memory incremental processing requires more than reducing recomputation or accelerating graph access alone: topology maintenance, dependency repair, and CPU--GPU data movement must preserve update locality throughout the processing pipeline.

% Challenge 1: preserving topological locality
First, local topology changes can incur nonlocal maintenance costs. An edge update directly modifies only its source's outgoing adjacency, but shared graph layouts may move unrelated adjacency lists during insertion or rebalancing. Such movement also invalidates their GPU adjacency descriptors, expanding both data movement and index maintenance beyond the updated sources. Exploiting topological locality therefore requires a layout that isolates adjacency relocation and a maintenance scheme that consistently coordinates outgoing and incoming adjacency views.

% Challenge 2: compact incremental computation through phase decoupling
Second, coupling deletion repair with insertion processing can make incremental computation unnecessarily broad. Without phase decoupling, the insertion phase must conservatively add all vertices to the active set to search widely for vertices whose results may need repair, incurring substantial computation as well as large CPU--GPU data-transfer overhead. We decouple the two phases so that, during deletion repair, the CPU supplies reliable replacement predecessors and obtains a stable result before insertion edges are processed. Consequently, both deletion repair and subsequent insertion propagation need to update only results related to the modified topology, greatly reducing the work range, computation, and CPU--GPU transfer time.

% Challenge 3: workload-dependent overhead
Third, update locality alone does not eliminate workload-dependent overhead. In large-diameter graphs, successive improvements may propagate over many rounds and repeatedly scan downstream adjacency. For large update batches, repeated preparation, copying, and sorting across maintenance stages can become a substantial CPU cost. These workloads require complementary optimizations that reduce redundant propagation and reuse batch-level maintenance information.

% Overview of CoIncGraph
To address these challenges, we present CoIncGraph, a CPU--GPU cooperative system for incremental processing of out-of-memory streaming graphs. CoIncGraph combines three complementary designs. Its \textit{source-isolated graph organization} stores each source's adjacency in an independently allocated block and coordinates topology changes through a unified change view, confining adjacency relocation and GPU index maintenance to changed sources. Its \textit{CPU--GPU cooperative compact incremental computation} exploits the locality of dynamic graph updates by decoupling insertion and deletion processing. With CPU assistance, it confines computation to affected vertices and their dependencies, thereby reducing both the computational workload and data transfer overhead of incremental processing. Finally, its \textit{workload-specific optimizations} use ordered propagation to reduce redundant edge scans and shared batch preparation to avoid repeated CPU work.

% Contributions
Our contributions are as follows:

1) We design a source-isolated graph organization that preserves topological locality and coordinates consistent maintenance of host adjacency and GPU indices.

2) We develop CPU--GPU cooperative compact incremental computation that exploits dynamic-update locality to limit the work set to affected vertices and their dependencies.

3) We introduce ordered propagation and shared batch preparation to reduce redundant work in large-diameter graphs and large update batches.

4) We evaluate CoIncGraph on six real-world graphs against Ingress, SHGraph, and Grapin, demonstrating its performance benefits across CPU-based, GPU in-memory, and CPU--GPU out-of-memory baselines.
